% Copyright (c) 2022-2026 MKM Research Labs. Licensed under the MIT License; see LICENSE.
%
% MKM-SS-001 Storm Sequence Generator -- model documentation v3.0, the document of record
% filed on the model as Core Documentation. It supersedes v2.0 (11-Mar-2026). The Wolfram
% benchmark of v2.0 (MKM-SS-001_wolfram_benchmark v0.1) found that the code computed its
% equations, but that the duration cap left no room for gaps in 46% of multi-storm
% sequences, the stated rates and Table 9 were wrong, the override control described as
% absent existed and reached the generator only in part, and the cited paths and evidence
% were stale. v3.0 reserves each gap's minimum before capping the storms, moves the tables
% to config so an override reaches the generator, and states every figure as computed.
%
% Build: it inputs the house header, which lives in PhysicalRisk. Copy
% PhysicalRisk/docs/models/shared beside this folder as ../shared, then run
% pdflatex three times (the third settles "Page n of N").
\documentclass[11pt,a4paper]{article}
\newcommand{\doctitle}{Storm Sequence Generator}
\newcommand{\docsubtitle}{Multi-Storm Sequence Model v3.0 (MKM-SS-001)}
\newcommand{\docversion}{3.0}
\newcommand{\docdate}{08-November-2026}
\newcommand{\docauthor}{David K Kelly}
\input{../shared/mkm_header}
% T1 so the underscores in code identifiers print as underscores.
\usepackage[T1]{fontenc}
\usepackage{lmodern}
% The shared history table sets every column to its natural width; wrap the description.
\renewcommand{\mkmhistorypage}[1]{%
    \section{Document History}
    \begin{table}[ht]
        \centering
        \begin{tabular}{l|l|p{9cm}|l}
            \textbf{Date} & \textbf{Version} & \textbf{Description} & \textbf{Author} \\
            \hline
            #1
        \end{tabular}
    \end{table}
}

\begin{document}
\sloppy
\mkmtitlepage
\mkmlegalpage
\tableofcontents
\clearpage

% ===================================================================
\mkmexecsummary
% ===================================================================

The Storm Sequence Generator draws the synthetic storm catalogue for the multi-storm
stress test: by default 10{,}000 \texttt{StormSequence}s, each holding 1--5 storms within a
168-hour event window. All rain ends by hour 156, which leaves at least 12 hours to watch
the rivers drain. The window follows the hours clause in catastrophe reinsurance. The
stress pipeline turns each sequence into a hydrograph at every gauge. Those hydrographs
train the per-gauge flood classifiers (MKM-FC-001) and drive the Portfolio Stress scenarios.

There are four sequence types:
\begin{itemize}[nosep]
    \item \textbf{isolated:} one storm;
    \item \textbf{doublet:} two storms with a medium gap of 24--72\,h;
    \item \textbf{cluster:} three or four storms with medium gaps;
    \item \textbf{persistent:} four or five storms with short gaps of 6--36\,h.
\end{itemize}
The chance of a multi-storm sequence rises with the intensity category, from 5\% for
minimal storms to 85\% for catastrophic ones.

\textbf{v3.0 keeps a gap between every pair of storms.} v2.0 capped each storm at
$156/n$\,h before drawing the gaps. Long storms then filled the window and left no time
between them: in 46\% of multi-storm sequences the storms ran back to back, so the
catchment recovery the gaps exist to model did not happen. v3.0 reserves each gap's
minimum (6\,h short, 24\,h medium) before it caps the storms, and shortens a long gap
only toward that minimum. Every multi-storm sequence now has a gap of at least its
minimum, and every sequence still fits the window without retries.

\textbf{v3.0 also corrects what v2.0 said.} The Wolfram benchmark of v2.0 (v0.1) found
that the code computed v2.0's equations but that 11 of its 17 claims did not hold:
\begin{itemize}[nosep]
    \item its cap and scaling rates and its Table~9 were wrong;
    \item its intensity correlation was misdescribed;
    \item its base precipitations were for catchments the platform does not have;
    \item it said no runtime override existed, but one does, and it reached the generator
          only in part;
    \item its paths, command line and classifier evidence were stale.
\end{itemize}
v3.0 moves every table to \texttt{config.port}, where an override reaches the generator.
Every figure here is generated from the code. Wolfram benchmark v0.2 tests each claim
(validation question Q8).

\textbf{Key risk: the durations do not fit the window.} Table~3's storms are too long for
three to five of them, with their gaps, to fit in 156\,h. The cap binds on 60.5\% of all
duration draws (\S6.2), and in 63.4\% of multi-storm sequences every storm is capped. This
exceeds the 50\% level at which this document asks for recalibration (Q6). Recalibrating
Tables~1 and~3 is outstanding (Q9). The classifiers on file were trained on v2.0's
sequences and must be retrained (\S7.3).

% ===================================================================
\mkmpurpose
% ===================================================================

\subsection{Purpose}
\begin{enumerate}[nosep]
    \item \textbf{Synthetic sequences.} A catalogue of storm sequences across the moderate
          to catastrophic categories, each with 1--5 timed storms.
    \item \textbf{Hydrograph timing.} Each storm's start, duration, intensity and peak
          position, so the stress pipeline can compose multi-peak hydrographs that show a
          catchment recovering between storms.
    \item \textbf{Classifier training data.} Hydrographs with antecedent wetness: a gauge
          that has only partly recovered from one storm is at greater risk from the next.
\end{enumerate}

\subsection{Where it is used}
\begin{itemize}[nosep]
    \item \texttt{app/commands/port/stages/storm.py}, the \texttt{-{}-stressm} stage of
          \texttt{phys.py port}, calls the stress pipeline.
    \item \texttt{src/port/src/stressm/pipeline/orchestrator/\_core.py} generates the
          sequences, saves them, and computes each gauge's response.
    \item \texttt{src/port/src/stressm/batch\_train.py} and
          \texttt{src/routes/trading/stress/training.py} load the saved sequences to train
          the classifiers (MKM-FC-001).
    \item The Portfolio Stress screens read the sequences and the classifiers' outputs.
\end{itemize}

\subsection{Scope}
In scope: sequence type, storm count, durations, gaps, intensities, precipitation, the
window, and batch generation. Out of scope: the spatial footprint (MKM-SG-001), the
gauge response and hydrology in the package's \texttt{models/} subpackage, hazard curves
(MKM-GH-001), pricing (MKM-PR-001), climate projections and tides.

\subsection{Tier}
Tier 2. The generator produces no regulated output itself, but it is the main input to the
stress classifiers, which feed the Portfolio Stress P\&L.

% ===================================================================
\mkmframework
% ===================================================================

\subsection{Event Window}
The event window $W$ and the hour $W_P$ by which all precipitation must end are
\begin{align}
    W &= 168\text{ h}, \label{eq:w}\\
    W_P &= W - W_{\text{drain}} = 168 - 12 = 156\text{ h}. \label{eq:wp}
\end{align}
A sequence's own drainage window runs from the end of its last storm to the end of the
event window:
\begin{equation}
    W^{(\text{seq})}_{\text{drain}} = W - (t_{\text{start},n} + d_n) \;\ge\; W_{\text{drain}} = 12\text{ h}.
    \label{eq:drain}
\end{equation}

\subsection{Sequence Types}
The storm count for a cluster or persistent sequence is drawn uniformly from its range.
\begin{table}[H]
\centering\small
\caption{Sequence types.}\label{tab:types}
\begin{tabular}{llll}
\toprule
Type & Storms & Gap type & Physical analogue \\
\midrule
isolated & 1 & --- & a single frontal system \\
doublet & 2 & medium & an Atlantic low followed by a secondary low \\
cluster & 3--4 & medium & an active frontal passage \\
persistent & 4--5 & short & a blocking pattern; near-continuous rain \\
\bottomrule
\end{tabular}
\end{table}

\subsection{Sequence Type Selection}
A storm becomes a multi-storm sequence with a probability that depends on its category:
\begin{equation}
P(\text{sequence}) = \begin{cases}
0.05 & \text{minimal} \\ 0.15 & \text{baseline} \\ 0.30 & \text{moderate} \\
0.50 & \text{severe} \\ 0.70 & \text{extreme} \\ 0.85 & \text{catastrophic}
\end{cases} \label{eq:pseq}
\end{equation}
If it does, the type is drawn from Table~\ref{tab:typew}.
\begin{table}[H]
\centering\small
\caption{Type weights for a multi-storm sequence, by category.}\label{tab:typew}
\begin{tabular}{lrrr}
\toprule
Category & Doublet & Cluster & Persistent \\
\midrule
minimal, baseline & 0.60 & 0.30 & 0.10 \\
moderate & 0.50 & 0.35 & 0.15 \\
severe & 0.30 & 0.50 & 0.20 \\
extreme & 0.20 & 0.45 & 0.35 \\
catastrophic & 0.15 & 0.45 & 0.40 \\
\bottomrule
\end{tabular}
\end{table}

\subsection{Storm Durations}
Each storm's duration is a triangular draw on its category's parameters
(Table~\ref{tab:dur}), rounded to the nearest whole hour:
\begin{equation}
    D = \operatorname{round}\big(\text{Triangular}(a_{\text{cat}}, c_{\text{cat}}, b_{\text{cat}})\big).
    \label{eq:dur}
\end{equation}
\begin{table}[H]
\centering\small
\caption{Storm duration (hours), triangular, by category.}\label{tab:dur}
\begin{tabular}{lrrr}
\toprule
Category & Min & Mode & Max \\
\midrule
minimal & 2 & 6 & 12 \\ baseline & 6 & 12 & 24 \\ moderate & 12 & 30 & 60 \\
severe & 24 & 48 & 96 \\ extreme & 36 & 72 & 144 \\ catastrophic & 48 & 120 & 240 \\
\bottomrule
\end{tabular}
\end{table}

\subsubsection{Duration Cap}
Each of the $n-1$ gaps is guaranteed at least its gap type's minimum $g_{\min}$ (6\,h short,
24\,h medium; Table~\ref{tab:gap}). The window reserves that time first, and divides what
is left equally between the storms:
\begin{equation}
    d_i = \min\!\left(D_i,\; \frac{W_P - (n-1)\,g_{\min}}{n}\right), \qquad i = 1, \dots, n.
    \label{eq:cap}
\end{equation}
The cap is 156\,h for an isolated storm, 66\,h in a doublet, 36\,h and 21\,h in clusters of
three and four, and 34.5\,h and 26.4\,h in persistent sequences of four and five. If the
window cannot hold $n$ storms and their minimum gaps (cap $\le 0$, possible only under an
override), generation stops with an error. A capped duration is not rounded.

\subsection{Inter-Storm Gaps}
Gaps are triangular draws on their gap type's parameters, rounded to the nearest whole hour.
\begin{table}[H]
\centering\small
\caption{Inter-storm gap (hours), triangular, by gap type.}\label{tab:gap}
\begin{tabular}{llrrrl}
\toprule
Gap type & Used by & Min & Mode & Max & Interpretation \\
\midrule
short & persistent & 6 & 18 & 36 & minimal drainage; high compounding \\
medium & doublet, cluster & 24 & 48 & 72 & partial drainage (30--60\% recovery) \\
long & --- & 48 & 96 & 144 & substantial drainage; unused \\
\bottomrule
\end{tabular}
\end{table}

\subsubsection{Gap Shortening}
Let $R = W_P - \sum_i d_i$ be the time the storms leave. Eq.~\eqref{eq:cap} guarantees
$R \ge (n-1)\,g_{\min}$. If the raw gaps $G_1, \dots, G_{n-1}$ overflow it,
\begin{equation}
    \sum_{j=1}^{n-1} G_j > R, \label{eq:overflow}
\end{equation}
each gap keeps its minimum and the time above the minimum is scaled to fit:
\begin{equation}
    g_j = g_{\min} + (G_j - g_{\min})\,\frac{R - (n-1)\,g_{\min}}{\sum_{k=1}^{n-1} (G_k - g_{\min})},
    \qquad j = 1, \dots, n-1. \label{eq:shorten}
\end{equation}
Otherwise the raw gaps are used. Either way $g_j \ge g_{\min}$ and
$\sum d_i + \sum g_j \le W_P$. When every storm is capped, $R = (n-1)\,g_{\min}$ and every
gap is held at its minimum.

\subsection{Storm Placement}
\begin{align}
    t_{\text{start},1} &= 0, \label{eq:t1}\\
    t_{\text{start},i} &= t_{\text{start},i-1} + d_{i-1} + g_{i-1}, \qquad i = 2, \dots, n. \label{eq:ti}
\end{align}

\subsection{Intensities}
The sequence's base intensity is drawn per category (Table~\ref{tab:base}) and floored:
\begin{equation}
    F_0 = \max\!\big(0.1,\; \mathcal{N}(\mu_{\text{cat}}, \sigma^2_{\text{cat}})\big). \label{eq:f0}
\end{equation}
\begin{table}[H]
\centering\small
\caption{Base intensity, by category.}\label{tab:base}
\begin{tabular}{lrr}
\toprule
Category & $\mu$ & $\sigma$ \\
\midrule
minimal & 0.30 & 0.10 \\ baseline & 0.60 & 0.15 \\ moderate & 1.00 & 0.20 \\
severe & 1.80 & 0.30 \\ extreme & 3.00 & 0.50 \\ catastrophic & 5.00 & 0.80 \\
\bottomrule
\end{tabular}
\end{table}
Each storm varies around $F_0$ by $\delta$, within $\pm v = \pm 0.20$. The first storm
varies downward:
\begin{equation}
    F^{(1)} = \max\!\big(0.1,\; F_0(1 + \delta_1)\big), \qquad \delta_1 \sim \text{Uniform}(-0.2, 0).
    \label{eq:first}
\end{equation}
The code draws $\delta_1$ and then, with probability 0.70, draws it again from the same law.
The redraw cannot change $\delta_1$'s law. Its configuration parameter,
\texttt{FIRST\_STORM\_DOMINANT\_PROB} (0.30, shown on the Trading Control screen), therefore
has no effect on the output distribution. It is kept only so the random stream is
unchanged. Each later storm varies up with probability $p_{\text{corr}} = 0.70$:
\begin{equation}
    F^{(i)} = \max\!\big(0.1,\; F_0(1 + \delta_i)\big), \qquad
    \delta_i \sim \begin{cases} \text{Uniform}(0, 0.2) & \text{with probability } 0.70, \\
    \text{Uniform}(-0.2, 0) & \text{with probability } 0.30. \end{cases} \label{eq:later}
\end{equation}
A later storm is then at least as strong as the first with probability
$p_{\text{corr}} + (1 - p_{\text{corr}})/2 = 0.85$ (above the floor): it is if it varies up,
and half the time if both vary down. The model represents antecedent wetness this way: a
catchment wet from one storm responds more to the next.

\subsection{Precipitation}
\begin{align}
    P_i &= P_{\text{base}}\, F^{(i)}, \label{eq:precip}\\
    \text{peak position} &\sim \text{Beta}(2, 2), \label{eq:peak}
\end{align}
where $P_{\text{base}}$ is the catchment's base precipitation (Table~\ref{tab:pbase}) and the
peak position is the fraction of the storm's duration at which it peaks.

\subsection{Aggregates}
\begin{align}
    P_{\text{total}} &= \textstyle\sum_i P_i, \label{eq:ptot}\\
    F_{\max} &= \max_i F^{(i)}, \label{eq:fmax}\\
    \bar F &= \tfrac{1}{n}\textstyle\sum_i F^{(i)}, \label{eq:fbar}\\
    F_{\text{cum}} &= \textstyle\sum_i F^{(i)}, \label{eq:fcum}\\
    D_{\text{total}} &= \textstyle\sum_i d_i + \sum_j g_j. \label{eq:dtot}
\end{align}
$F_{\text{cum}}$ carries the compounding of several storms into the stress test, where the
peak alone would not.

% ===================================================================
\mkminputs
% ===================================================================

\subsection{Batch Generation}
\texttt{generate\_event\_set(count, catchment\_id, intensity\_weights, force\_sequence\_type,
seed)} draws each sequence's category from Table~\ref{tab:weights} (normalised, so custom
weights need not sum to one). The default count is 10{,}000.
\begin{table}[H]
\centering\small
\caption{Default category weights.}\label{tab:weights}
\begin{tabular}{lr}
\toprule
Category & Weight \\
\midrule
moderate & 0.40 \\ severe & 0.35 \\ extreme & 0.20 \\ catastrophic & 0.05 \\
\bottomrule
\end{tabular}
\end{table}

\subsection{Catchment Base Precipitation}
$P_{\text{base}}$ is \texttt{BASE\_PRECIPITATION\_MM} in the catchment's storm module,
\texttt{catch/<catchment>/storm.py}. Three catchments have one (Table~\ref{tab:pbase}). A
catchment without one (the single-module catchments under \texttt{catch/}, such as the
Rhine) takes the Thames's 35\,mm.
\begin{table}[H]
\centering\small
\caption{Base precipitation by catchment.}\label{tab:pbase}
\begin{tabular}{lr}
\toprule
Catchment & $P_{\text{base}}$ (mm) \\
\midrule
thames & 35.0 \\ mekong & 90.0 \\ halong & 100.0 \\
any other & 35.0 \\
\bottomrule
\end{tabular}
\end{table}

\subsection{Constants}
\begin{table}[H]
\centering\small
\caption{Constants.}\label{tab:const}
\begin{tabular}{llrl}
\toprule
Constant & Symbol & Value & Meaning \\
\midrule
Event window & $W$ & 168\,h & the hours clause \\
Minimum drainage & $W_{\text{drain}}$ & 12\,h & watched after the last storm \\
Precipitation window & $W_P$ & 156\,h & $W - W_{\text{drain}}$ \\
Intensity variation & $v$ & 0.20 & each storm's range around $F_0$ \\
Correlation probability & $p_{\text{corr}}$ & 0.70 & $P(\delta_i \ge 0)$ for a later storm \\
Later storm $\ge$ first & & 0.85 & $p_{\text{corr}} + (1-p_{\text{corr}})/2$ \\
Peak shape & $(\alpha,\beta)$ & (2, 2) & peak position \\
Intensity floor & & 0.1 & on $F_0$ and every $F^{(i)}$ \\
\bottomrule
\end{tabular}
\end{table}

\subsection{Where the Parameters Live, and Overrides}
Every parameter in Tables~2--8, except $P_{\text{base}}$, is in \texttt{config/port/\_storm.py}
(\texttt{config.port}). The generators read them from there on every draw.
\texttt{storm\_control.json}, in the catchment's input directory, overrides any of them:
\begin{itemize}[nosep]
    \item The Trading Control screen writes the file and applies it when it saves.
    \item Importing \texttt{config} applies it too, so a command-line run uses it.
\end{itemize}
Because the generators read \texttt{config.port} on every draw, an override takes effect
whether it was applied before or after the generator was imported. A file that sets some
rows of a table changes those rows and keeps the others. Saving on the screen requires the
trading write permission. \texttt{force\_sequence\_type} forces every sequence's type, for
testing.

% ===================================================================
\mkmimplementation
% ===================================================================

\subsection{Software Architecture}
The package is \texttt{src/port/src/storm\_multi/}:
\begin{itemize}[nosep]
    \item \texttt{core/data\_structures.py}: \texttt{SequenceStorm}, \texttt{StormSequence},
          \texttt{SequenceType}, \texttt{GapType}, \texttt{make\_storm\_id},
          \texttt{make\_sequence\_id}.
    \item \texttt{generators/sequence\_generator.py}: \texttt{SequenceGenerator}, which
          applies Eqs.~\eqref{eq:cap}--\eqref{eq:dtot}.
    \item \texttt{generators/batch\_generator.py}: \texttt{generate\_event\_set}.
    \item \texttt{generators/duration\_sampler.py}: \texttt{sample\_duration},
          \texttt{get\_duration\_range}.
    \item \texttt{generators/gap\_sampler.py}: \texttt{sample\_gap}, \texttt{get\_gap\_range},
          \texttt{get\_min\_gap}.
    \item \texttt{generators/intensity\_sampler.py}: \texttt{should\_generate\_sequence},
          \texttt{sample\_sequence\_type}, \texttt{get\_storm\_count},
          \texttt{sample\_storm\_intensity}.
    \item \texttt{utils/validation.py}: \texttt{validate\_sequence},
          \texttt{validate\_event\_set}, \texttt{fits\_in\_window}, \texttt{precip\_end\_hour},
          \texttt{min\_gap\_hours}.
    \item \texttt{utils/serialization.py}: \texttt{save\_sequences}, \texttt{save\_summary},
          \texttt{load\_sequences}.
\end{itemize}
The \texttt{models/} subpackage (sequence response, hydrology, spatial correlation) is part
of the gauge response, not of this model.

\subsection{Generation Flow}
For each of the $N$ sequences, \texttt{generate\_event\_set}:
\begin{enumerate}[nosep]
    \item draws the category (Table~\ref{tab:weights}) and $F_0$ (Eq.~\eqref{eq:f0});
    \item draws the type (Eq.~\eqref{eq:pseq}, Table~\ref{tab:typew}) and the storm count;
    \item draws each storm's duration and caps it (Eqs.~\eqref{eq:dur}--\eqref{eq:cap}), and
          its intensity (Eqs.~\eqref{eq:first}--\eqref{eq:later});
    \item draws the gaps and shortens them if they overflow
          (Eqs.~\eqref{eq:overflow}--\eqref{eq:shorten});
    \item places the storms (Eqs.~\eqref{eq:t1}--\eqref{eq:ti}), draws each peak position,
          and computes the precipitation and aggregates.
\end{enumerate}

\subsection{Reproducibility}
Two \texttt{numpy.random.RandomState} streams are seeded from the one seed: one draws the
categories and base intensities, the other everything in a sequence. A fixed seed
reproduces the batch exactly, apart from the identifiers and timestamps. Each sequence also
carries a per-event seed for the storm--typhoon coupling, from a separate stream that does
not disturb the storm draws.

\subsection{Validation}
\texttt{validate\_event\_set} checks every sequence:
\begin{itemize}[nosep]
    \item the type is one of the four, and the storm count is 1--5 and equals the storms held;
    \item there are $n-1$ gaps, each at least its type's minimum;
    \item each storm starts after the previous one ends;
    \item the last storm ends by $W_P$, leaving at least $W_{\text{drain}}$;
    \item every intensity category is valid;
    \item the total duration matches the durations and gaps.
\end{itemize}
Times are compared to within $10^{-9}$\,h. Eqs.~\eqref{eq:cap} and~\eqref{eq:shorten} make
the timing checks hold by construction; validation is a second line of defence.

\subsection{Command Line and Outputs}
\begin{verbatim}
python phys.py port --stressm [--num-sims N] [--gauge-id GID] [--seed N]
\end{verbatim}
\texttt{-{}-num-sims} sets the count (default 10{,}000). The stage deletes the catchment's
classifiers, which were trained on the previous catalogue, then runs the stress pipeline.
That pipeline saves \texttt{storm\_sequences.json} (the catalogue) and
\texttt{sequences\_summary.json} (batch statistics) for the catchment. The classifiers are
then retrained per gauge from the Portfolio Stress screen.

% ===================================================================
\mkmsensitivity
% ===================================================================

Every figure in this section is exact. It is computed in rational arithmetic from the
tables above, with durations and gaps rounded as the code rounds them.

\subsection{Sequence Types}
For the default weights (Table~\ref{tab:weights}), a sequence is isolated with
probability $\sum_c w_c\,(1 - P_c(\text{sequence}))$. A multi-storm type has probability
$\sum_c w_c\,P_c(\text{sequence})\,\pi_{c,t}$, where $\pi_{c,t}$ is its weight in
Table~\ref{tab:typew}.
\begin{table}[H]
\centering\small
\caption{Expected sequence types in a default batch of 10{,}000.}\label{tab:t9}
\begin{tabular}{lrr}
\toprule
Type & Expected count & Share \\
\midrule
isolated & 5{,}225 & 52.25\% \\
doublet & 1{,}469 & 14.69\% \\
cluster & 2{,}116 & 21.16\% \\
persistent & 1{,}190 & 11.90\% \\
\bottomrule
\end{tabular}
\end{table}
The doublet and cluster counts are 1{,}468.75 and 2{,}116.25 exactly. By storm count, a batch
of 10{,}000 holds 5{,}225 sequences of one storm, 1{,}468.75 of two, 1{,}058.13 of three,
1{,}653.13 of four and 595 of five.

\subsection{Duration Cap}
The cap of Eq.~\eqref{eq:cap} binds on \textbf{60.5\%} of all duration draws in the default
batch. It binds on every catastrophic duration in a five-storm sequence: the cap there is
26.4\,h, and a catastrophic storm lasts at least 48\,h. Table~\ref{tab:cap} gives each
configuration.
\begin{table}[H]
\centering\small
\caption{Share of duration draws capped (\%), by configuration and category.}\label{tab:cap}
\begin{tabular}{lrrrrrr}
\toprule
Type & $n$ & Cap (h) & Moderate & Severe & Extreme & Catastrophic \\
\midrule
isolated & 1 & 156.0 & 0.0 & 0.0 & 0.0 & 30.3 \\
doublet & 2 & 66.0 & 0.0 & 25.2 & 76.1 & 97.5 \\
cluster & 3 & 36.0 & 38.4 & 91.0 & 100.0 & 100.0 \\
cluster & 4 & 21.0 & 89.6 & 100.0 & 100.0 & 100.0 \\
persistent & 4 & 34.5 & 45.2 & 93.6 & 100.0 & 100.0 \\
persistent & 5 & 26.4 & 75.7 & 99.6 & 100.0 & 100.0 \\
\bottomrule
\end{tabular}
\end{table}

\subsection{Gap Shortening}
Gaps are shortened (Eq.~\eqref{eq:shorten}) in \textbf{81.1\%} of multi-storm sequences.
That is 99.0\% of moderate-to-severe persistent sequences and 100.0\% of catastrophic ones
(all but two in $10^{10}$). By type, the rate is 39.5\% of doublets, 99.6\% of clusters
and 99.5\% of persistent sequences. In \textbf{63.4\%} of multi-storm sequences every storm
is capped, which holds every gap at its minimum (Table~\ref{tab:gaps}). A shortened gap is
never shorter than its minimum.
\begin{table}[H]
\centering\small
\caption{Multi-storm sequences whose gaps are shortened, and (in brackets) whose every
storm is capped (\%), by configuration and category. 100.0 is above 99.95.}\label{tab:gaps}
\begin{tabular}{lrrrrr}
\toprule
Type & $n$ & Moderate & Severe & Extreme & Catastrophic \\
\midrule
doublet & 2 & 0.9 (0.0) & 47.9 (7.2) & 92.9 (60.2) & 99.7 (95.6) \\
cluster & 3 & 96.4 (7.2) & 100.0 (78.8) & 100.0 (100.0) & 100.0 (100.0) \\
cluster & 4 & 100.0 (70.5) & 100.0 (100.0) & 100.0 (100.0) & 100.0 (100.0) \\
persistent & 4 & 93.9 (4.2) & 100.0 (76.8) & 100.0 (100.0) & 100.0 (100.0) \\
persistent & 5 & 100.0 (24.8) & 100.0 (98.2) & 100.0 (100.0) & 100.0 (100.0) \\
\bottomrule
\end{tabular}
\end{table}
Most multi-storm sequences therefore take their gaps from the minimum of Table~\ref{tab:gap}
rather than its full range. The gaps themselves are 6--36\,h and 24--72\,h only where the
storms are short.

% ===================================================================
\mkmvalidation
% ===================================================================

\subsection{Structural Validation}
\texttt{validate\_event\_set} finds no invalid sequence in a default batch, or in 400
sequences of every category with every forced type. Each sequence ends by hour 156, and
every gap is at least its minimum.

\subsection{Type Distribution}
Pooled over 100 seeds of 10{,}000 sequences, every type's count is within 1\% of
Table~\ref{tab:t9}. A single batch varies more: its persistent count has a standard error
of about 3\%.

\subsection{Downstream Classifier}
The classifiers' training summary on file (\texttt{training\_summary.json}, 31 October
2026) covers 11 Thames gauges, with a mean AUC-ROC of 0.9935. Those classifiers were
trained on v2.0's sequences, and so is the catalogue on file. Both predate v3.0, and its
longer gaps change the hydrographs. The next \texttt{-{}-stressm} run regenerates the
catalogue and deletes the classifiers, which must then be retrained. MKM-FC-001 owns that
evidence. This document does not claim the classifiers' performance for v3.0's sequences.

\subsection{An Example Gauge}
GAUGE-1aaf8b08, in the same summary: AUC-ROC 0.9929, accuracy 0.9776, Brier score 0.0169;
flood rate 8.6\% over 1{,}713{,}600 training vectors. The dominant features are
\texttt{log\_h\_s} (76\%) and \texttt{log\_t\_end} (17\%).

% ===================================================================
\mkmlimitations
% ===================================================================

\begin{enumerate}[nosep]
    \item \textbf{The durations do not fit the window.} With Table~\ref{tab:dur}'s
          durations, three to five severe or worse storms and their minimum gaps cannot
          fit in 156\,h. The cap binds on 60.5\% of draws and truncates the most extreme
          multi-storm sequences, understating their duration. Recalibrating Tables~1
          and~3 (fewer storms for long categories, or shorter storms in sequences) is
          outstanding.
    \item \textbf{Gaps are mostly at their minimum.} In 63.4\% of multi-storm sequences
          every gap is held at $g_{\min}$, so recovery between storms is the least the
          gap type allows rather than a draw from its range.
    \item \textbf{Gap shortening is proportional, not physical.} The time above each
          gap's minimum is scaled uniformly.
    \item \textbf{Storms are independent within a sequence.} Durations and peak positions
          carry no serial correlation, and a long first storm does not shorten the second.
    \item \textbf{One footprint.} Every storm in a sequence has the same spatial footprint.
    \item \textbf{$P_{\text{base}}$ is fixed per catchment.} It does not vary by season or
          wetness, and a catchment without a storm module takes the Thames's value.
    \item \textbf{The season is the generation date.} Each sequence's
          \texttt{sequence\_start} is the time it was generated. The stress pipeline reads
          that month as the season for each gauge's base level, so every sequence in a
          batch uses the base levels of the month the batch was run.
    \item \textbf{A control with no effect.} \texttt{FIRST\_STORM\_DOMINANT\_PROB} does not
          change the output distribution (\S3.7).
    \item \textbf{Stationarity.} No climate scenario changes the weights, persistence or
          intensities.
\end{enumerate}

% ===================================================================
\section{Validation Questions}
% ===================================================================

\mkmvqitem{1}{Purpose}{Is the model being used for its intended purpose?}
Yes. It generates the stress catalogue that the stress pipeline turns into hydrographs for
the classifiers and the Portfolio Stress scenarios (\S2.2).

\mkmvqitem{2}{Conceptual Soundness}{Is the model conceptually sound?}
The structure is sound: typed sequences in a fixed window, a guaranteed gap between storms,
and correlated intensities. Its parameters are not mutually consistent, however, because
the storms are too long for the window (\S6.2, \S8 item~1).

\mkmvqitem{3}{Data Quality}{Is the data quality adequate?}
The parameters are expert-set, not fitted to observed sequences. No observed sequence
catalogue has been compared with them.

\mkmvqitem{4}{Developmental Evidence}{Is there adequate developmental evidence?}
Unit tests cover the samplers, the cap and gap minimum, the validator, the override path
and reproducibility.

\mkmvqitem{5}{Outcomes Analysis}{Are outcomes reasonable?}
The sequences are valid by construction and the type mix is as stated. Whether sequences
like these occur over the Thames is untested.

\mkmvqitem{6}{Ongoing Monitoring}{What ongoing monitoring is required?}
On each regenerated catalogue, monitor:
\begin{itemize}[nosep]
    \item the type counts against Table~\ref{tab:t9};
    \item the cap rate: a rate above 50\% calls for recalibration of the durations, and the
          baseline rate is 60.5\%;
    \item the gap-shortening rate by type: a rate above 70\% says the gaps and durations
          are inconsistent with the window, and the baseline is 99.6\% for clusters and
          99.5\% for persistent sequences;
    \item each retrained classifier's AUC-ROC, flagging a gauge below 0.95.
\end{itemize}
The cap and shortening thresholds are breached at the baseline, as \S8 records.

\mkmvqitem{7}{Benchmarking}{Has the model been benchmarked?}
Not against observed sequences or another generator. The implementation is benchmarked
against Wolfram (Q8).

\mkmvqitem{8}{Implementation}{Is the implementation verified?}
Wolfram oracle benchmark v0.2 tests every equation and claim in this document. It replays
19{,}600 sequences number for number, and checks the exact probabilities in Wolfram.

\mkmvqitem{9}{Overall Risk}{Is the overall model risk acceptable?}
Acceptable for its Tier 2 use, with two outstanding items: the durations must be
recalibrated so that multi-storm sequences are not mostly capped (\S8 item~1), and the
classifiers must be retrained on v3.0's catalogue (\S7.3). Until then, the stress
scenarios understate the duration of the most extreme sequences.

\mkmhistorypage{11-Mar-2026 & 2.0 & First issue: the multi-storm sequence generator & David K Kelly \\
08-Nov-2026 & 3.0 & Gap minima reserved before the duration cap, so every multi-storm
sequence keeps its gaps; tables moved to config so an override reaches the generator;
figures, correlation, base precipitation, paths and evidence stated as computed (Wolfram
benchmark v0.1) & David K Kelly \\}

\end{document}
